\honest{Extensions and Intensions}{Eliezer Yudkowsky}{2008}

``What is red?'' ``Red is a color.'' ``What's a color?'' ``A color is a property of a thing.'' Soon both speakers are lost in words defined by other words, the problem Steven Harnad compared to learning Chinese from a Chinese/Chinese dictionary. Instead I could point to a stop sign, a red shirt, blood. I think I stole this example from S. I. Hayakawa, ``though I'm really not sure.'' \nb{The memory is right. Hayakawa's \textsc{Language in Thought and Action} has the same exchange about ``red'' and the same answer: look at a traffic light.}

To define a word by other words, as a dictionary does, is to give an ``intensional definition''; to point to examples is to give an ``extensional definition.'' \nb{These are Hayakawa's senses. In logic, an extensional definition lists every member of a class, and pointing to a few examples is called an ostensive definition.}

Rationalists in films are lost in words, but Charles Sanders Peirce, a real one, defined lithium by a recipe: search among certain minerals, treat and fuse and dissolve them, and you get ``a globule of a pinkish silvery metal that will float on gasolene.'' This, I say, is a ``treasure map'' that leads to an example. \nb{Peirce called it ``a precept that is more serviceable than a definition.''}

Now the concepts themselves. The actual intension of my ``tiger'' concept is ``the neural pattern (in my temporal cortex)'' that decides whether what I see is a tiger; its actual extension is ``everything I call a tiger.'' \nb{In logic, an intension is the property a word connotes, not a pattern in anyone's brain. The post uses the word in its own sense.} Neither words nor pointing conveys the whole concept. If I point at one tiger, you may think I mean ``yellow thing.'' \nb{Wittgenstein made this point about ostensive definition in 1953.} So definitions are ``treasure maps, not treasure,'' and you cannot program your concepts into someone else's brain.

Nor can you choose both a definition and what it picks out. Whether ``a huge red rocky sphere'' matches the red light I call Mars is a separate matter, and my mind goes on believing that ``Mars'' is the fourth planet whatever definition I write. Because this matching happens below awareness, I conclude, arguing that something is true ``by definition'' is so popular. \nb{The post gives no example of such an argument; later posts in the sequence do.}
